\section{Introduction}\label{sec:introduction}

For $\abs{x}\le 1$ let
\[
  \chi_2(x)=\sum_{j\ge0}\frac{x^{2j+1}}{(2j+1)^2}=\tfrac12\bigl(\Li_2(x)-\Li_2(-x)\bigr),
  \qquad
  \Ti_2(x)=\sum_{j\ge0}\frac{(-1)^jx^{2j+1}}{(2j+1)^2}=\operatorname{Im}\Li_2(ix)
\]
be Legendre's chi function and the inverse tangent integral. This paper presents a proof of the
following statement.

\begin{theorem}\label{thm:main}
Let $a\in\{3,4,5,8\}$ and let $q$ be an odd integer prime to $a$.
\begin{enumerate}[(i)]
\item The number $\chi_2(a/q)$ is irrational if $q\ge q_+(a)$, where
$q_+(3)=131$, $q_+(4)=995$, $q_+(5)=829$ and $q_+(8)=12061$.
\item The number $\Ti_2(a/q)$ is irrational if $q\ge q_-(a)$, where
$q_-(3)=257$, $q_-(4)=995$, $q_-(5)=1609$ and $q_-(8)=12061$.
\end{enumerate}
Since $\chi_2$ and $\Ti_2$ are odd, the same holds for $\chi_2(-a/q)$ and $\Ti_2(-a/q)$.
\end{theorem}

The thresholds are inclusive, and they are sufficient for our estimates; we do not claim that
they are optimal (\cref{sec:limitations}).

\emph{Status.} The proof is a candidate. It was reconstructed and checked by AI systems from
candidate arguments produced by another AI system (see the disclosure at the end); no
independent human expert has reviewed it. In Lean only the final deduction is formalised, and
only conditionally on the three estimates of \cref{thm:three-estimates}. All finite
inequalities used in the proof have been recomputed by programs written separately from those
that first produced them (Appendices~\ref{app:certificates} and~\ref{app:consistency}); this
checks the numbers, not the asymptotic arguments. \Cref{subsec:step-status} gives the status of
each step, and \cref{rem:scrutiny} lists the points that most need independent scrutiny.

\subsection{Background}\label{subsec:background}

If $q$ is even or not prime to $a$, the fraction $a/q$ reduces to one with a smaller numerator
or an even numerator and odd denominator; we treat the four numerators $3,4,5,8$ with odd
denominators prime to them. The irrationality of $\chi_2$ and $\Ti_2$ at rational points with
large denominator follows from linear independence criteria for shifted polylogarithms. Hata's
measures~\cite{Hata1990} are stated at reciprocals of integers. The criterion of David,
Hirata-Kohno and Kawashima~\cite{DHK2022Acta} applies to the shifted dilogarithm
$\sum_{n\ge1}x^n/(n-\frac12)^2$ at rational points; Krisman evaluates it at $x=4/q^2$ and
obtains the irrationality of $\chi_2(2/q)$ for $q\ge11\,795$ \cite[\S1.2]{Krisman2026chi2}.
The same evaluation at $x=\pm a^2/q^2$ gives the range $q>54e^4a^2$, that is, $q$ larger than
about $2.7\cdot10^4$, $4.7\cdot10^4$, $7.4\cdot10^4$ and $1.9\cdot10^5$ for $a=3,4,5,8$. We have not re-derived this range from
\cite{DHK2022Acta}; we quote it only to indicate the order of magnitude of what was known, and
we have not found explicit smaller thresholds for these numerators in the literature.

For small denominators the closest works are Krisman's preprint~\cite{Krisman2026chi2}, which
proves that $\chi_2(1/q)\notin\Q$ for $q\ge3$, $q\ne4$, that $\Ti_2(1/q)\notin\Q$ for $q\ge5$,
and that $\chi_2(2/q)\notin\Q$ for odd $q\ge175$, and the AxiomForge
preprint~\cite{AxiomForge2026numeratorTwo}, a candidate proof that $\chi_2(2/q)$ and
$\Ti_2(2/q)$ are irrational for odd $q\ge81$. Both are unrefereed and AI-assisted. We use none
of their theorems as an input: every estimate needed here is proved in this paper, including
those that are restatements, for general numerator, of estimates
of~\cite{AxiomForge2026numeratorTwo}. The method is a Hankel-determinant construction in the
spirit of Zudilin's determinantal approach~\cite{Zudilin2017}: Krisman takes the Hankel
determinants of the functional $g\mapsto\sum_{t\ge1}w^t\,(g(u)\,u/D(u))'|_{u=t}$, where
$D(u)=\prod_{k<h}(u+k+\frac12)$ has as many half-integer poles as the size $h$ of the
determinant; in~\cite{AxiomForge2026numeratorTwo} the number of poles exceeds $h$ in the ratio
$7/6$.

\subsection{The construction}\label{subsec:constructions}

Fix an integer $a\ge1$, an odd integer $q>a$ prime to $a$ and a sign $\sigma\in\{1,-1\}$, and
put
\begin{equation}\label{eq:basic-notation}
  y=\frac qa,\qquad w=\frac{\sigma}{y^2}=\frac{\sigma a^2}{q^2},\qquad W=\abs w,\qquad
  \xi=\sum_{j\ge0}\frac{\sigma^jy^{-2j-1}}{(2j+1)^2},
\end{equation}
so that $\xi=\chi_2(a/q)$ for $\sigma=1$ and $\xi=\Ti_2(a/q)$ for $\sigma=-1$. For integers
$K\ge h\ge1$ let
\[
  x_k=-k-\tfrac12\ (0\le k<K),\qquad D_K(u)=\prod_{k<K}(u-x_k),\qquad
  \varphi(u)=\frac{u}{D_K(u)}.
\]
In \cref{sec:pencils} we define an affine functional $\mu_X$ on the rational functions with
simple poles at the $x_k$ such that $\mu_\xi(R)=\sum_{m\ge1}w^mR'(m)$, and the polynomial
\[
  T(X)=\nrm\cdot\det\bigl[\mu_X(\varphi(u)\,u^{i+j})\bigr]_{0\le i,j<h}\in\Q[X],
  \qquad \nrm=\frac{K!^{\,h}}{\prod_{i<h}i!^2},
\]
of degree at most $h$. We use two families of sizes, indexed by $n\ge1$:
\[
  (h,K)=(6n,7n)\quad(\kappa=\tfrac76),\qquad (h,K)=(4n,5n)\quad(\kappa=\tfrac54),
\]
where $\kappa=K/h$. The first is used for $\chi_2(3/q)$ and the second for the seven other
cases of \cref{thm:main}. Throughout,
\[
  \Dk=\frac{(2-\kappa)^2}{4},\qquad\text{so}\qquad D_{7/6}=\frac{25}{144},\quad
  D_{5/4}=\frac9{64}.
\]

\subsection{The argument in one page}\label{subsec:one-page}

For a prime $p$ and a nonzero $F\in\Q[X]$ let $\vG p(F)$ be the minimum of the $p$-adic
valuations of the coefficients of $F$. Let $\ell=\prod_pp^{-\vG p(T)}$, so that $\ell T\in\Z[X]$.
If $\xi=r/s$ were rational, then $s^h\ell T(r/s)$ would be an integer. The proof consists in
showing that this integer is nonzero for infinitely many $n$, while its absolute value is at
most $\exp\bigl(h^2(\bconst-\aconst+o(1))\bigr)$ with $\bconst<\aconst$
(\cref{prop:criterion}). This needs three estimates.

\begin{theorem}\label{thm:three-estimates}
Let $a\ge1$, let $q>a$ be odd and prime to $a$, let $\sigma=\pm1$ and $\kappa\in\{7/6,5/4\}$,
and let $T=T_n$ be the polynomial attached to the $n$-th size of ratio $\kappa$, with $h=h_n$.
Write $a=2^ea_0$ with $a_0$ odd.
\begin{enumerate}[(i)]
\item \emph{Arithmetic.} Along the $n$ with $T_n\ne0$,
\begin{multline*}
  \liminf\frac{1}{h^2}\sum_p\vG p(T)\log p\ \ge\ \aconst_\sigma(a,q)\\
  :=\Iprof_\kappa+c_2\log2-(2\kappa-1)\log a_0+\frac13\log q-\Dk\log(q^2-\sigma a^2),
\end{multline*}
where $\Iprof_{7/6}=-2.113164137$, $\Iprof_{5/4}=-2.355639833$, and
\[
  c_2=\begin{cases}2(\kappa-1)+2\Dk&\text{if $a$ is odd and }\sigma=1,\\
  \kappa-1&\text{if $a$ is odd and }\sigma=-1,\\
  C_e(\kappa)=\min_{0\le x\le1}f_e(x)&\text{if }e\ge1,\end{cases}
\]
with $f_e(x)=2(\kappa-1)-2(e+1)\kappa x+(e+2)x^2+(\frac e2+2)(1-x)^2$.
\item \emph{Size.} For every comparison density $\rho$ and every $\lambda$ satisfying the tail
condition of \cref{lem:tail}, $\limsup h^{-2}\log\abs{T(\xi)}\le\bconst(\rho,\lambda)$, an
explicit number that depends on $(a,q)$ only through $W$.
\item \emph{Nonvanishing.} Let $r\in\Q$. If $q\ge14a$, then $T_n(r)\ne0$ for every
sufficiently large prime $n$.
\end{enumerate}
\end{theorem}

\begin{proof}
Part~(i) is \cref{thm:arithmetic}, part~(ii) is \cref{thm:decay}, and part~(iii) is
\cref{thm:nonvanishing} together with \cref{prop:local-dets}, since $q\ge14a$ gives
$\abs{\sigma q^2/a^2}\ge196>189$.
\end{proof}

\Cref{tab:constants} lists the constants at the six thresholds. For $a=4$ and $a=8$ the two
signs are treated together, with the smaller arithmetic constant $\aconst_-$
(\cref{rem:worse-sign}).

\begin{table}[ht]
\centering
\caption{The six threshold comparisons, per $h^2$. The arithmetic constant is
$\aconst_\sigma(a,q_0)$ of \cref{thm:three-estimates}, with the value of $c_2$ shown; $\bconst$ is
the certified size exponent (Appendix~\ref{app:decay-certificates}). Decimals are rounded in the
safe direction.}
\label{tab:constants}
\small
\begin{tabular}{@{}llllllll@{}}
\toprule
$a$ & $q_0$ & function & $\kappa$ & $c_2$ & $\aconst\ge$ & $\bconst\le$ & $\aconst-\bconst\ge$\\
\midrule
$3$ & $131$ & $\chi_2$ & $7/6$ & $49/72$ & $-3.173875358$ & $-3.187821144$ & $0.013945786$\\
$3$ & $257$ & $\Ti_2$ & $5/4$ & $1/4$ & $-3.541276253$ & $-3.546389753$ & $0.005113500$\\
$4$ & $995$ & both & $5/4$ & $-337/112$ & $-4.081754611$ & $-4.083454335$ & $0.001699724$\\
$5$ & $829$ & $\chi_2$ & $5/4$ & $25/32$ & $-3.878258885$ & $-3.879559185$ & $0.001300300$\\
$5$ & $1609$ & $\Ti_2$ & $5/4$ & $1/4$ & $-4.211960841$ & $-4.212674322$ & $0.000713481$\\
$8$ & $12061$ & both & $5/4$ & $-9/2$ & $-4.985336979$ & $-4.986095509$ & $0.000758530$\\
\bottomrule
\end{tabular}
\end{table}

The estimates are obtained as follows.
\begin{itemize}
\item \emph{Prime $2$, $a$ odd} (\cref{subsec:two-odd}). Here $y$ is a $2$-adic unit, and $2$
divides both the denominators of the poles and $q^2-\sigma a^2$. In bases of polynomials that
vanish at the first and at the last poles, each entry of the matrix is a forward difference of a
sequence obeying a first-order recursion, and its valuation is bounded directly.
\item \emph{Prime $2$, $a$ even} (\cref{subsec:two-even}). In the coordinate $z=2u$ the poles
are odd and the summation points even, and a Vandermonde-type minor whose nodes share a parity
is divisible by about $2^{t^2}$; the Cauchy--Binet formula reduces $\vG2(T)$ to the minimisation
of $f_e$ in the proportion $x$ of pole terms.
\item \emph{Odd primes dividing $a$} (\cref{subsec:p-divides-a}). The pole $x_k$ carries the
factor $y^{2k+1}$, of valuation $-(2k+1)v_p(a)$, but it enters the matrix through a term of rank
one, so a term of the determinant uses each pole at most once; the $h$ most expensive poles cost
$(2\kappa-1)v_p(a)h^2$.
\item \emph{Odd primes dividing $q(q^2-\sigma a^2)$} (\cref{sec:exceptional}). For $p\mid q$,
reflecting $m\mapsto-m$ makes the series $p$-adically convergent, and each level offers at most
three terms of valuation at least $2v_p(q)m$, which gives $\frac13v_p(q)h^2$; for
$p\mid q^2-\sigma a^2$ the only loss comes from the denominators $(1-w)^{b+1}$, and a convexity
argument gives $\Dk\,v_p(q^2-\sigma a^2)h^2$.
\item \emph{Ordinary primes} (\cref{sec:ordinary}). For $p\asymp h$ two lower bounds, by Taylor
jets in each residue class modulo $p$ and by global partial fractions with Vandermonde gains,
are small minimisation problems over the residue classes; their maximum, integrated against
$\log p$ by the prime number theorem, gives $\Iprof_\kappa$, and the integral is evaluated
exactly (Appendix~\ref{app:profile-certificate}).
\item \emph{Size} (\cref{sec:decay}). The determinant at $X=\xi$ is a sum over configurations
of $h$ points on the summation sites; Vandermonde bounds and a discrete energy inequality
compare it with an explicit step density, and the resulting number $\bconst$ is evaluated in
ball arithmetic (Appendix~\ref{app:decay-certificates}).
\item \emph{Nonvanishing} (\cref{sec:nonvanishing}). For $h=tp$ with $p$ prime, a basis adapted
to the residue classes modulo $p$ makes the scaled determinant block triangular modulo $p$, with
diagonal blocks equal, up to units, to three fixed local Hankel determinants, which are
nonzero.
\item \emph{Assembly} (\cref{sec:proofs}). The margins of \cref{tab:constants} are positive,
and for $q\ge q_0\ge2a$ the margin grows by at least $\frac{10}{27}\log(q/q_0)$
(\cref{prop:propagation}), which carries each threshold to all larger admissible $q$.
\end{itemize}

\subsection{What is new and what is inherited}\label{subsec:new}

The functional, the criterion, the reflection at primes dividing $q$, the size bound with its
shift to smaller $\abs w$, and the strategy of proving nonvanishing modulo a prime proportional
to $h$ go back to Krisman~\cite{Krisman2026chi2}. The ratio $K/h>1$, the parity-minor bound at
$2$ for $a=2$, the residue-class profile for the ordinary primes and the Taylor-jet form of the
nonvanishing argument are those of~\cite{AxiomForge2026numeratorTwo}, where $\kappa=7/6$. We
reprove all of them for general $a$ and both ratios; in
\cref{sec:ordinary,sec:decay,sec:nonvanishing} the numerator enters only through $w$, and the
proofs are those of~\cite{AxiomForge2026numeratorTwo} with a simpler weight. The following is
new.
\begin{enumerate}[(a)]
\item \emph{Primes dividing the numerator} (\cref{prop:p-divides-a}), which do not occur for
$a\le2$ apart from the prime $2$.
\item \emph{The prime $2$ for odd $a$} (\cref{thm:two-odd}). Its cost depends on
$\nu=v_2(q^2-\sigma a^2)$, but the dependence cancels against the contribution of $2$ to
$\log(q^2-\sigma a^2)$, so that $\aconst_\sigma(a,q)$ contains the fixed multiples
$\frac{49}{72}\log2$ or $\frac{25}{32}\log2$ ($\sigma=1$) and $\frac16\log2$ or $\frac14\log2$
($\sigma=-1$).
\item \emph{The prime $2$ for $a=2^ea_0$ with $e\ge2$} (\cref{thm:two-even}): the parity-minor
argument of~\cite{AxiomForge2026numeratorTwo} with pole costs $e(2k+1)$ and site costs $2em$.
For $e=3$ and $\kappa=5/4$ the minimum of $f_e$ is attained at $x=1$ and equals $-e(2\kappa-1)$,
the value given by the rank-one argument of~(a).
\item \emph{The ratio $5/4$}: a second profile integral, a second triple of local determinants
(of size $4$), and five new comparison densities.
\item \emph{Propagation in $q$ for general $a$} (\cref{prop:propagation}).
\end{enumerate}

\subsection{Organisation and notation}\label{subsec:notation}

\Cref{sec:pencils} defines the functional, the determinants and the criterion.
\Cref{sec:prime-two,sec:exceptional,sec:ordinary} bound the denominators, \cref{sec:decay} the
size, and \cref{sec:nonvanishing} proves nonvanishing; \cref{sec:proofs} assembles the proof and
records the status of each step, and \cref{sec:limitations} discusses what the thresholds do and
do not say. Appendix~\ref{app:certificates} describes the finite computations and
Appendix~\ref{app:consistency} compares the estimates with exact computations at small sizes.
\Cref{tab:notation} lists the notation used throughout.

\begin{table}[ht]
\centering
\caption{Notation.}\label{tab:notation}
\small
\begin{tabular}{@{}p{0.27\textwidth}p{0.66\textwidth}@{}}
\toprule
$a$, $q$, $\sigma$; $y$, $w$, $W$, $\xi$ & numerator, odd denominator prime to $a$, sign $\pm1$;
see~\eqref{eq:basic-notation}\\
$e$, $a_0$; $Q$, $\nu$ & $a=2^ea_0$ with $a_0$ odd; $Q=\sigma y^2=1/w$, $\nu=v_2(Q-1)$ for odd
$a$\\
$\tau_k$ & partial sum $\sum_{j\le k}\sigma^jy^{-2j-1}/(2j+1)^2$ of $\xi$, \eqref{eq:tau}\\
$h$, $K$; $\kappa$, $\Dk$ & size, number of poles; $K/h\in\{\frac76,\frac54\}$,
$(2-\kappa)^2/4$\\
$x_k$, $D_K$, $\varphi$, $\Phi$ & $-k-\frac12$; $\prod_{k<K}(u-x_k)$; $u/D_K$; $K!\,\varphi$\\
$\Omega_b$, $\mom$, $\mu_X$ & $w^{\max(b,1)}(1-w)^{-b-1}$; rational moment sum; the functional
\eqref{eq:functional}\\
$\Delta$, $\nrm$, $T$ & Hankel determinant, normaliser, $T=\nrm\Delta$, see~\eqref{eq:Delta-T}\\
$\beta_i$, $\vvec(u)$, $\bvec(u)$ & $\binom ui$; the columns $(u^i)_{i<h}$ and
$(\beta_i(u))_{i<h}$\\
$\vG p$ & Gauss valuation of a polynomial in $X$ over $\Q_p$\\
$\aconst_\sigma(a,q)$, $\bconst$, $\Iprof_\kappa$ & arithmetic constant, size exponent,
ordinary-prime constant\\
$\psi(t)$, $\pos x$, $s_2(n)$ & $t\log\abs t-t$ with $\psi(0)=0$; $\max(x,0)$; binary digit
sum, so $v_2(n!)=n-s_2(n)$\\
\bottomrule
\end{tabular}
\end{table}

\emph{Conventions.} $p$ is always a prime and $v_p$ the $p$-adic valuation on $\Q_p$. The Gauss
valuation satisfies $\vG p(FG)=\vG p(F)+\vG p(G)$ and $\vG p(F+G)\ge\min(\vG p(F),\vG p(G))$, with
$\vG p(0)=+\infty$, and is invariant under $X\mapsto X+c$ for $c\in\Z_p$; hence the determinant of
a matrix over $\Q_p[X]$ has Gauss valuation at least the minimum over the terms of its Leibniz
expansion. For $u$ in any $\Q$-algebra, $\binom uk=u(u-1)\cdots(u-k+1)/k!$. The implied constants
in $O(\cdot)$ may depend on $a$, $q$ and $\kappa$, but not on $n$.
